কতগুলি বস্তুর সমষ্টিকে \emph{সেট} বলে। ওই বস্তুগুলিই সেটটির
উপাদান। $x$ যদি~$A$-এর একটি উপাদান হয়, তবে আমরা $x \in A$ লিখি।
সেট \emph{সদস্যভিত্তিক সমতার নীতি} মেনে চলে। তার উপাদানগুলিই তাকে
সম্পূর্ণ নির্ধারণ করে। উপাদানগুলির স্পষ্ট \emph{তালিকা} দিয়ে অথবা
তাদের একটি সাধারণ ধর্ম উল্লেখ করে (\emph{ধর্মনির্দেশে সেট-গঠন}) সেট
নির্দিষ্ট করা যায়। সদস্যভিত্তিক সমতার নীতি অনুযায়ী, কোন ক্রমে বা
কীভাবে সেটের উপাদানগুলি তালিকায় লেখা বা নির্দিষ্ট করা হয়, তাতে কিছু
যায় আসে না। $A$ ও~$B$ একই সেট ($A = B$) প্রমাণ করতে হলে দেখাতে হবে
যে~$A$-এর প্রত্যেক উপাদান~$B$-এরও উপাদান, এবং বিপরীত দিক থেকেও একই
কথা সত্য।

কয়েকটি গুরুত্বপূর্ণ সেট হলো স্বাভাবিক সংখ্যার সেট ($\Nat$),
পূর্ণসংখ্যার সেট ($\Int$), মূলদ সংখ্যার সেট ($\Rat$) ও বাস্তব
সংখ্যার সেট ($\Real$)। বস্তু দিয়ে গঠিত \emph{প্রতীকক্রমের} ($X^*$)
ও অসীম \emph{অনুক্রমের} ($X^\omega$) সেটও গুরুত্বপূর্ণ। $A$-এর প্রত্যেক উপাদান~$B$-এরও
উপাদান হলে~$A$ হলো~$B$-এর একটি \emph{উপসেট}, $A \subseteq B$।
একটি সেট~$B$-এর সব উপসেটের সমষ্টি নিজেই একটি সেট, যার নাম~$B$-এর
\emph{ঘাত সেট} $\Pow{B}$। সেটের \emph{সংযোগ} $A \cup B$ এবং
\emph{ছেদ} $A \cap B$ গঠন করা যায়। একটি \emph{ক্রমযুগল}
$\tuple{x, y}$ দুটি বস্তু $x$ ও~$y$ নিয়ে গঠিত, ঠিক ওই ক্রমে।
$x = y$ না হলে $\tuple{x, y}$ ও $\tuple{y, x}$ ভিন্ন যুগল।
$x \in A$ ও $y \in B$ হয় এমন সব যুগল $\tuple{x, y}$-এর সেটকে
$A$ ও $B$-এর \emph{কার্তেসীয় গুণফল}~$A \times B$ বলে।
$A \times A$-কে আমরা $A^2$ লিখি। যেমন, $\Nat^2$ হলো স্বাভাবিক
সংখ্যার যুগলগুলির সেট।
